% fig03_architecture_timeline.tex — 图3 架构演进时间线
% 《深度学习与大语言模型：前沿技术全景报告》图示库
% 由 dl_llm_report.tex 抽取，经 % fig03_architecture_timeline.tex — 图3 架构演进时间线
% 《深度学习与大语言模型：前沿技术全景报告》图示库
% 由 dl_llm_report.tex 抽取，经 % fig03_architecture_timeline.tex — 图3 架构演进时间线
% 《深度学习与大语言模型：前沿技术全景报告》图示库
% 由 dl_llm_report.tex 抽取，经 % fig03_architecture_timeline.tex — 图3 架构演进时间线
% 《深度学习与大语言模型：前沿技术全景报告》图示库
% 由 dl_llm_report.tex 抽取，经 \input{figs/fig03_architecture_timeline} 引用（caption/label 在主文件）
% 注意：本文件为片段，依赖主文件导言区的 tikz/pgfplots 宏包、颜色定义与 tikzset 样式，不可单独编译。

\begin{tikzpicture}[x=1.25cm]
\draw[arr] (-0.4,0) -- (14.5,0) node[lab, below left=0mm and -2mm] {年份};
\foreach \i/\yr/\side/\sty/\txt in {
  0/1998/1/gry/{LeNet\\CNN 时代开启},
  1/2012/-1/gry/{AlexNet\\深度学习爆发},
  2/2014/1/gry/{Seq2Seq+注意力\\RNN 时代雏形},
  3/2015/-1/gry/{ResNet\\残差连接},
  4/2017/1/blk/{\textbf{Transformer}\\范式革命},
  5/2018/-1/blk/{BERT / GPT-1\\预训练+微调},
  6/2020/1/blk/{GPT-3\\上下文学习涌现},
  7/2022/-1/orn/{ChatGPT\\对齐走向产品},
  8/2023/1/blk/{LLaMA / GPT-4\\开源与多模态},
  9/2024/-1/orn/{o1 / DeepSeek-V2\\推理模型元年},
  10/2025/1/orn/{R1 / Qwen3 / GPT-5\\开源推理普惠},
  11/2026/-1/orn/{Agentic·世界模型\\从对话到行动}} {
  \node[\sty, font=\scriptsize, align=center] at (\i*1.25, \side*1.5) {\txt};
  \draw[cgray, thin] (\i*1.25,0) -- (\i*1.25, \side*0.55);
  \node[lab, font=\scriptsize] at (\i*1.25, -\side*0.4) {\yr};
}
\end{tikzpicture}
 引用（caption/label 在主文件）
% 注意：本文件为片段，依赖主文件导言区的 tikz/pgfplots 宏包、颜色定义与 tikzset 样式，不可单独编译。

\begin{tikzpicture}[x=1.25cm]
\draw[arr] (-0.4,0) -- (14.5,0) node[lab, below left=0mm and -2mm] {年份};
\foreach \i/\yr/\side/\sty/\txt in {
  0/1998/1/gry/{LeNet\\CNN 时代开启},
  1/2012/-1/gry/{AlexNet\\深度学习爆发},
  2/2014/1/gry/{Seq2Seq+注意力\\RNN 时代雏形},
  3/2015/-1/gry/{ResNet\\残差连接},
  4/2017/1/blk/{\textbf{Transformer}\\范式革命},
  5/2018/-1/blk/{BERT / GPT-1\\预训练+微调},
  6/2020/1/blk/{GPT-3\\上下文学习涌现},
  7/2022/-1/orn/{ChatGPT\\对齐走向产品},
  8/2023/1/blk/{LLaMA / GPT-4\\开源与多模态},
  9/2024/-1/orn/{o1 / DeepSeek-V2\\推理模型元年},
  10/2025/1/orn/{R1 / Qwen3 / GPT-5\\开源推理普惠},
  11/2026/-1/orn/{Agentic·世界模型\\从对话到行动}} {
  \node[\sty, font=\scriptsize, align=center] at (\i*1.25, \side*1.5) {\txt};
  \draw[cgray, thin] (\i*1.25,0) -- (\i*1.25, \side*0.55);
  \node[lab, font=\scriptsize] at (\i*1.25, -\side*0.4) {\yr};
}
\end{tikzpicture}
 引用（caption/label 在主文件）
% 注意：本文件为片段，依赖主文件导言区的 tikz/pgfplots 宏包、颜色定义与 tikzset 样式，不可单独编译。

\begin{tikzpicture}[x=1.25cm]
\draw[arr] (-0.4,0) -- (14.5,0) node[lab, below left=0mm and -2mm] {年份};
\foreach \i/\yr/\side/\sty/\txt in {
  0/1998/1/gry/{LeNet\\CNN 时代开启},
  1/2012/-1/gry/{AlexNet\\深度学习爆发},
  2/2014/1/gry/{Seq2Seq+注意力\\RNN 时代雏形},
  3/2015/-1/gry/{ResNet\\残差连接},
  4/2017/1/blk/{\textbf{Transformer}\\范式革命},
  5/2018/-1/blk/{BERT / GPT-1\\预训练+微调},
  6/2020/1/blk/{GPT-3\\上下文学习涌现},
  7/2022/-1/orn/{ChatGPT\\对齐走向产品},
  8/2023/1/blk/{LLaMA / GPT-4\\开源与多模态},
  9/2024/-1/orn/{o1 / DeepSeek-V2\\推理模型元年},
  10/2025/1/orn/{R1 / Qwen3 / GPT-5\\开源推理普惠},
  11/2026/-1/orn/{Agentic·世界模型\\从对话到行动}} {
  \node[\sty, font=\scriptsize, align=center] at (\i*1.25, \side*1.5) {\txt};
  \draw[cgray, thin] (\i*1.25,0) -- (\i*1.25, \side*0.55);
  \node[lab, font=\scriptsize] at (\i*1.25, -\side*0.4) {\yr};
}
\end{tikzpicture}
 引用（caption/label 在主文件）
% 注意：本文件为片段，依赖主文件导言区的 tikz/pgfplots 宏包、颜色定义与 tikzset 样式，不可单独编译。

\begin{tikzpicture}[x=1.25cm]
\draw[arr] (-0.4,0) -- (14.5,0) node[lab, below left=0mm and -2mm] {年份};
\foreach \i/\yr/\side/\sty/\txt in {
  0/1998/1/gry/{LeNet\\CNN 时代开启},
  1/2012/-1/gry/{AlexNet\\深度学习爆发},
  2/2014/1/gry/{Seq2Seq+注意力\\RNN 时代雏形},
  3/2015/-1/gry/{ResNet\\残差连接},
  4/2017/1/blk/{\textbf{Transformer}\\范式革命},
  5/2018/-1/blk/{BERT / GPT-1\\预训练+微调},
  6/2020/1/blk/{GPT-3\\上下文学习涌现},
  7/2022/-1/orn/{ChatGPT\\对齐走向产品},
  8/2023/1/blk/{LLaMA / GPT-4\\开源与多模态},
  9/2024/-1/orn/{o1 / DeepSeek-V2\\推理模型元年},
  10/2025/1/orn/{R1 / Qwen3 / GPT-5\\开源推理普惠},
  11/2026/-1/orn/{Agentic·世界模型\\从对话到行动}} {
  \node[\sty, font=\scriptsize, align=center] at (\i*1.25, \side*1.5) {\txt};
  \draw[cgray, thin] (\i*1.25,0) -- (\i*1.25, \side*0.55);
  \node[lab, font=\scriptsize] at (\i*1.25, -\side*0.4) {\yr};
}
\end{tikzpicture}
